% Background --- section structure from prd.md 3.1 (Ingenieur, Chapter 2), ownership from Table 3.
%
% This chapter OWNS the foundations of the engineering document: speech pipelines, multirotor
% vehicles and their control, swarm control, real-time systems and latency budgets, the
% single-board computer as a shared host, state machines and command buses.
% It does NOT own the foundations the Memoire de Master defines (language models, LoRA,
% quantisation formats, CPU inference and thermal throttling, GBNF, intent/slot SLU, metrics):
% those get one sentence each in sec:bg-master, never a re-explanation (write-once).
% No results, no design decisions, no positioning: what THIS system chose belongs in
% chap:architecture / chap:implementation, what the LITERATURE measured in chap:state-of-the-art.

\chapter{Background}
\label{chap:background}
This chapter sets out the foundations the rest of the document assumes, for a reader who has not
worked in speech processing, aerial robotics or real-time computing. It follows the path of a spoken
command. It begins with the stages that turn a stream of audio into words or detections, then
describes the vehicles those detections command and the control of a group of them. It then covers
the timing discipline a system of this kind is held to, the processor it runs on, and the two
software structures that carry a command to the vehicles: a state machine and a command bus. The
foundations of the language model itself are defined in the \emph{M\'emoire de Master} and are named
in Section~\ref{sec:bg-master} without being repeated. Every section states established material.
What published systems have built from these elements is reviewed in
Chapter~\ref{chap:state-of-the-art}, and the decisions made in this work are stated and justified
from Chapter~\ref{chap:architecture} onwards.

\section{Speech pipelines}
\label{sec:bg-speech}
A speech pipeline turns a continuous audio signal into discrete decisions. The signal reaches the
software as a sequence of samples, commonly 16{,}000 per second for speech~\cite{speechcmd,whisper},
and is processed in \emph{frames}: short, fixed-length blocks of consecutive samples, tens of
milliseconds long. The stages of a pipeline are of two kinds. Some run on every frame, continuously,
and decide something about a short window of audio. Others run once per utterance and need the whole
of it before they can begin. Stages of the first kind determine how soon a system can react; stages
of the second kind, which see the whole utterance, determine what it can interpret.

\paragraph{Keyword spotting.} A keyword spotter detects a small, fixed set of phrases in a
continuous audio stream. It is a task distinct from full speech recognition, with its own datasets
and evaluation methodology~\cite{speechcmd}. Its output is a detection score per frame rather than a
transcript, and a detection is declared when that score crosses a threshold. A spotter is small
enough to run on every frame, so it can listen continuously while larger models stay idle. Its
errors are of two kinds. A \emph{false accept} is a detection when the phrase was not spoken, and is
counted per hour of audio that does not contain it. A \emph{false reject} is a spoken phrase that
was not detected, and is counted as a proportion of the times it was spoken. One utterance can hold
the score above the threshold for several consecutive frames, so a spotter commonly \emph{debounces}
its detections: it ignores a frame above the threshold if another crossed it within a short
\emph{refractory} interval before~\cite{speechcmd}.

\paragraph{Operating point.} A single threshold trades the two errors against each other. Raising
it removes false accepts and adds false rejects; lowering it does the opposite. The curve traced by
the two error rates as the threshold varies is a \gls{roc} curve, and the threshold at which a
spotter is operated is its \emph{operating point}. A false-accept rate quoted without the
false-reject rate at the same threshold is therefore uninformative, since a rate as low as desired
can be reached by raising the threshold far enough. The threshold is chosen on
data held apart from the data on which the resulting error rates are reported, or the reported
rates are optimistic.

\paragraph{Voice activity detection.} A \gls{vad} model classifies each short window of audio as
speech or non-speech. The simplest detectors threshold the energy of the signal; trained detectors
are small neural classifiers that output a probability of speech per window and remain usable in
noise that has the energy of speech but not its structure. A \gls{vad} model does not recognise
words. Its role in a pipeline is to decide which stretches of audio are worth passing to the more
expensive stages downstream.

\paragraph{Endpointing.} Endpointing is the decision that an utterance has ended. It is commonly
made from the \gls{vad} output by a silence rule: the utterance is declared over once the detector
has reported non-speech for a set duration. That duration sets a trade-off. A short silence
threshold answers sooner, but it cuts an utterance at a pause between words and passes an incomplete
command downstream. A long one keeps pauses inside the utterance, but adds its full length to the
delay before every command. The delay is a wait rather than a computation: no processor speed
shortens it. In this document, the \emph{end of speech} is the instant at which the endpointer
declares the utterance over. It follows the last spoken word by about the silence threshold, and it
is the earliest instant at which a stage that needs the complete utterance can begin.

\paragraph{Speech recognition.} \Gls{stt}, or speech recognition, maps the audio of an utterance to
a sequence of words. Many current recognisers are neural sequence-to-sequence models: an encoder
turns the audio into a sequence of representations, and a decoder generates the transcript from them
one token at a time. The Whisper models of the original release are transformers of this kind,
trained on 680{,}000 hours of weakly supervised audio in sizes from 39~million parameters
upwards~\cite{whisper}. Within one family, a larger model is more accurate and slower, and because
the decoder is autoregressive, the time taken also grows with the length of the transcript.
Recognition accuracy is scored by the \gls{wer}, a metric the \emph{M\'emoire de Master} defines
(Section~\ref{sec:bg-master}).

\paragraph{Noise and signal-to-noise ratio.} The \gls{snr} of a recording compares the power of
the speech with the power of the noise added to it, on a logarithmic scale:
\begin{equation}
  \label{eq:snr}
  \mathrm{SNR} = 10 \log_{10} \frac{P_{\mathrm{speech}}}{P_{\mathrm{noise}}} \quad \text{dB},
\end{equation}
where each power is the mean of the squared samples. At 0~dB the two powers are equal, and every
10~dB fall multiplies the relative power of the noise by ten. A controlled robustness test mixes a
recorded noise into clean speech, scaling the noise so that the mixture reaches a chosen \gls{snr},
and measures how recognition degrades as that ratio falls.

\paragraph{Cascaded and direct pipelines.} A spoken command can be interpreted in two ways. A
\emph{cascade} first transcribes the utterance and then passes the transcript to a separate stage
that extracts its meaning; errors made by the recogniser reach the second stage as text, without the
audio that could have resolved them. A \emph{direct} pipeline maps audio to a decision without
producing a transcript, as a keyword spotter does for its fixed phrases. The cascade can interpret
any sentence its recogniser can transcribe. The direct pipeline is limited to the decisions it was
trained on, and it avoids the cost of transcription. A direct pipeline that decides on every frame,
as a keyword spotter does, also avoids the endpointing wait; one that classifies a complete
utterance must wait for its end, as a cascade does. The meaning a cascade extracts from a
transcript, an intent and its slot values, is the frame defined in the \emph{M\'emoire de Master}
(Section~\ref{sec:bg-master}).

\section{Multirotor vehicles and flight control}
\label{sec:bg-uav}
A \gls{uav} is an aircraft flown without a pilot on board, either by a remote operator or by
software. The vehicles considered in this document are multirotors, and specifically quadrotors,
which lift and steer themselves with four fixed-pitch rotors. The combined thrust of the rotors
opposes gravity, and differences in thrust and torque between them tilt and turn the airframe.
A quadrotor cannot push itself sideways directly. To move horizontally it tilts, so that part of
its thrust points in the direction of travel, and to stop it tilts the other way. Hovering is the
state in which the thrust balances the weight and the airframe is level. It is not a state of rest
for the motors, which keep turning to hold the vehicle up. \TODO{a multirotor-dynamics reference,
e.g.\ Mahony, Kumar and Corke, IEEE Robotics \& Automation Magazine 2012 --- not in
references.bib; add and verify, then cite here}

\paragraph{Energy of a moving vehicle.} A flying vehicle carries kinetic energy in its velocity and
potential energy in its height. A command that stops horizontal motion removes kinetic energy, and
one that brings the vehicle down removes potential energy; a command that starts motion or climbs
adds energy. The distinction matters for safety: the harm of a collision grows with the kinetic
energy at impact, and that of a fall with the height it starts from.

\paragraph{The flight controller.} Flight control is organised as nested loops. The innermost loops,
on the flight controller carried by the vehicle, regulate the angular rate and the attitude of the
airframe at several hundred cycles per second. Outer loops regulate velocity and position more
slowly, by requesting an attitude and a thrust from the inner ones. Software commanding a vehicle
from outside therefore does not drive the rotors. It sends a \emph{setpoint} --- a target position,
velocity or acceleration --- to the outermost loop that accepts one, and the flight controller
tracks it. Open-source flight-control software exposes such setpoint interfaces, and so do
simulators that model the flight controller~\cite{pyflyt}. \TODO{a flight-stack source for the
nested loops, their rates and the setpoint interfaces, e.g.\ the PX4 Autopilot user guide, pages
Controller Diagrams and multicopter filter tuning --- not in references.bib; add and verify, then
cite here}

\paragraph{PID control.} The loops of a flight controller, and many controllers built above them,
use the \gls{pid} law. For an error $e(t)$ between a setpoint and the measured value, the command is
\begin{equation}
  \label{eq:pid}
  u(t) = k_p\, e(t) + k_i \int_0^t e(\tau)\, d\tau + k_d\, \frac{de(t)}{dt},
\end{equation}
where the three gains weight three corrections. The proportional term pulls towards the setpoint in
proportion to the error. The integral term removes an error that persists, such as the droop
caused by a constant disturbance. The derivative term opposes the rate at which the error changes,
and so damps the approach. When the output saturates, the integral keeps accumulating error the
actuator cannot act on, and the controlled value overshoots while that surplus unwinds after the
saturation ends. This \emph{integral wind-up} is limited by bounding the integral, among other
methods. \TODO{a feedback-control textbook for the PID law and
critical damping, e.g.\ Astrom and Murray, \emph{Feedback Systems} --- not in references.bib}

\paragraph{The double integrator.} The simplest model of a vehicle treats it as a point mass whose
acceleration is commanded directly, so that velocity is the integral of acceleration and position
the integral of velocity. Under a proportional and derivative law on the position error, this model
responds like a mass on a spring with a damper. With unit mass, the response is \emph{critically
damped} when $k_d = 2\sqrt{k_p}$: starting from rest, it reaches the setpoint in the shortest time
that any $k_d$ allows for that $k_p$ without overshooting. A smaller $k_d$ overshoots and oscillates
about the setpoint, and a larger one approaches more slowly. In discrete time the integration is
done in steps of a fixed duration. The \emph{semi-implicit} step updates the velocity first and
moves the position with the new velocity. For an undamped oscillator and a small enough step it
keeps the amplitude bounded, whereas updating both from the old values makes the amplitude grow at
every step.

\section{Swarm control}
\label{sec:bg-swarm}
In this document, a swarm is a group of vehicles commanded as one. Two families of methods
coordinate such a group, and they answer different requests. Flocking produces coordinated motion
without prescribing where any individual vehicle goes. Formation control assigns each vehicle a
place in a shape and drives it there. Formation control needs a separate mechanism to keep the
vehicles apart while they move; flocking contains one among its rules.

\paragraph{Flocking.} Reynolds' model of flocking produces coordinated group motion from three rules
applied by each individual to its nearby flockmates: avoid collisions with them, match their
velocity, and stay close to them~\cite{reynolds1987}. The model needs no leader and no global plan,
and the group behaviour emerges from rules that are local to each individual. Applied to the whole
group rather than to each individual's neighbours, the centring rule becomes a pull towards the
group's centroid, which Reynolds calls a \emph{central-force} model. He replaced it with his local
rule because it draws every member of a scattered flock towards the centroid at
once~\cite{reynolds1987}. Velocity matching applied to the whole group becomes, in the same way, a
pull towards the group's mean velocity.

\paragraph{Formations and slot assignment.} A formation is a set of positions, or \emph{slots},
defined relative to a reference point and a heading: equal angles on a circle, even spacing on a
line, two trailing arms for a wedge. Commanding a formation involves two steps. The first assigns
each vehicle to a slot. The second drives each vehicle to its slot, for instance with a \gls{pid}
law on the vehicle's error to it. The assignment decides how far each vehicle travels and whether
two of them can meet on the way. Assigning each vehicle to its nearest free slot in turn can send
two vehicles through the same place at the same time. Minimising the total squared distance over all
vehicles at once, an instance of the linear assignment problem, prevents such a meeting under two
conditions. The vehicles move in straight lines and arrive together, and their starting positions,
like their slots, are spaced widely enough for their size. Two paths may then still cross, but the
vehicles reach the crossing at different times. \TODO{Turpin, Michael and Kumar, ``CAPT: Concurrent
assignment and planning of trajectories for multiple robots'', International Journal of Robotics
Research 33(1), 2014 --- not in references.bib; add and verify, then cite here}

\paragraph{Separation and potential fields.} Separation is the requirement that no two vehicles
come closer than a set distance. An \emph{artificial potential field} treats obstacles, including
other vehicles, as sources of a repulsive force that grows as the distance to them shrinks, and
the goal as a source of attraction. A common repulsive potential between two vehicles at distance
$d$ is
\begin{equation}
  \label{eq:apf}
  U(d) = \tfrac{1}{2}\,\eta \left(\frac{1}{d} - \frac{1}{d_0}\right)^{2} \ \text{for } d \leq d_0,
  \qquad U(d) = 0 \ \text{otherwise},
\end{equation}
where $\eta$ sets its strength and $d_0$ the distance beyond which it has no effect. The magnitude
of the resulting force is $\eta\,(1/d - 1/d_0)/d^2$, which grows without bound as $d$ approaches
zero. \TODO{Khatib, the original artificial-potential-field formulation (ICRA 1985, or the journal
version of 1986; Koren and Borenstein cite the 1985 paper) --- not in references.bib; add and
verify, then cite eq:apf} Potential fields have failure modes identified as
inherent to the method: trap situations at local minima where attraction and repulsion cancel,
oscillation near obstacles and in narrow passages, and no passage between closely spaced
obstacles~\cite{koren1991}. A force also acts only through the dynamics. A closing pair stops short
of the separation distance only if the potential there exceeds the pair's relative kinetic energy,
and with bounded acceleration and a discrete timestep even contact becomes reachable. A potential
field is therefore a soft mechanism, which shapes motion, as opposed to a hard one, which
constrains the state directly.

\paragraph{Simulation.} Controllers for aerial vehicles are developed and tested in
\emph{software-in-the-loop} simulation: the control software runs unchanged against a simulated
vehicle in place of a real one. Simulators differ in fidelity. A kinematic simulator integrates the
point-mass model of Section~\ref{sec:bg-uav} and runs many times faster than real time. A physics
simulator integrates rigid-body and rotor dynamics, usually at a finer timestep than the controller
runs at, and models the vehicle's own attitude control. PyFlyt, for instance, provides \gls{uav}
simulation environments for reinforcement-learning research on the Bullet physics
engine~\cite{pyflyt}. Environments of this kind expose two operations: one that places the vehicles
at the start of a trial, and one that applies a command and advances the simulation by one step.
A trial seeded with a fixed number starts from the same state every time it is run, which makes a
simulated experiment exactly repeatable.

\section{Real-time systems and latency budgets}
\label{sec:bg-realtime}
A real-time system is one whose correctness depends on when its results are produced as well as on
what they are. A result delivered after its deadline is a failure even if its content is right. In
a \emph{hard} real-time system a missed deadline is a failure of the system, even if the late result
is right. In a \emph{soft} one it degrades the service, and requirements state how rarely a deadline
may be missed rather than promising that it never is. \TODO{a real-time systems textbook for the
hard/soft distinction and preemptive scheduling, e.g.\ Buttazzo, \emph{Hard Real-Time Computing
Systems} --- not in references.bib}

\paragraph{Latency as a distribution.} The latency of an operation is the time between two defined
events, and it is fixed only once both events are named: a figure measured from the start of an
utterance and one measured from its end describe different quantities. Repeated runs of the same
operation take different times, because what else the processor is doing, what its caches hold and
how hot it is vary from run to run. Latency is therefore reported as percentiles of a sample of
runs. By the nearest-rank method, the $q$-th percentile of $n$ sorted measurements is the one at
rank $\lceil q n / 100 \rceil$. The median, p50, is the typical run; p95 and p99 describe the tail,
since about one run in twenty is slower than p95 and one in a hundred slower than p99. A requirement
written on the mean says nothing about the tail. A requirement on p95 bounds the slowest runs a user
will regularly meet while tolerating the occasional exception.

\paragraph{Latency budgets.} A latency budget states a target for an end-to-end latency and an
allowance for each stage it passes through. Its purpose is attribution. When the end-to-end target
is missed, a stage that exceeded its own allowance is the first place to look. Allowances that
bound worst cases can be added, because the worst case of a sum is at most the sum of the worst
cases. Percentiles do not add: the
p95 of a sum of stage times is not the sum of the stages' p95 values, because the slow runs of
different stages need not coincide. An end-to-end percentile target is therefore judged on the
end-to-end measurement itself, and the stage allowances serve as ceilings for diagnosis, not as
shares that must sum to the total.

\paragraph{Periodic work.} A stage that processes a stream, such as the frame loop of a speech
pipeline, is a periodic task. It receives a new frame every period and must finish its work on each
frame before the next one arrives. If it overruns, the following frames wait in a queue and their
results are late until the backlog drains; if it overruns on average, the queue grows without
bound. The constraint on
such a task is on the worst frame, not on the average one.

\paragraph{Preemption and cancellation.} \emph{Preemption} is the interruption of running work in
favour of more urgent work. An operating-system scheduler preempts a running thread to give its
processor to another. An application can also withdraw work it has started: a long computation
checks a cancellation flag at safe points, such as between the tokens of a generated sequence, and
stops at the next check once the flag is set. Withdrawn work is discarded rather than resumed, and
this document calls its withdrawal preemption as well. The latency of such a cancellation is bounded
by the longest interval between checks, plus the time needed to release what the computation held.
It helps the urgent work only if that work can run while the work to be withdrawn still holds the
processor, for instance on a core of its own.

\section{The single-board computer as a host}
\label{sec:bg-sbc}
The \emph{M\'emoire de Master} describes the \gls{sbc} of this work, the Raspberry~Pi~5, and
explains why decoding a language model on its \gls{cpu} is bounded by memory bandwidth and why its
clock falls when it overheats. This section adds what that description does not need: how several
programs with different deadlines share the board's four cores.

\paragraph{A general-purpose operating system.} The board runs Linux, a general-purpose operating
system whose scheduler shares the processor cores among all runnable threads. Under its default
policy it aims at throughput and fairness, and it gives no timing guarantee to any one thread: a
thread that is ready to run can wait while others hold the cores. Linux also provides real-time
scheduling policies, which give a thread strict priority over ordinary threads. The operating system
does allow a program to restrict a process or a thread to a subset of the cores, a setting called
\emph{processor affinity}. A thread restricted to one core, while every other thread is restricted
to the rest, has that core to itself, apart from kernel work and interrupts, which affinity set by a
program does not move. \TODO{a source for Linux scheduling policies and processor affinity, e.g.\
the sched(7) and sched\_setaffinity(2) manual pages --- not in references.bib}

\paragraph{Shared resources.} Affinity separates computation, not everything the cores share. The
cores of a system-on-chip draw on the same memory bandwidth and, depending on the design, the same
last-level cache, and they heat the same die. Work on one core can therefore slow work on another
without ever being scheduled there, by saturating memory accesses or by raising the temperature
until the clock falls. On a board with four cores, a program with a tight deadline and a program
that saturates memory bandwidth interfere through these resources however their threads are
placed.

\section{State machines and command buses}
\label{sec:bg-fsm}
Two software structures carry a command from the stage that produces it to the stage that acts on
it: a command bus, which transports it, and a state machine, which decides whether it may take
effect.

\paragraph{Finite-state machines.} A finite-state machine has a finite set of states, a set of
inputs, and a transition function that gives, for each state and input, the next state and any
action taken. An input that has no transition from the current state is rejected in that state.
The legality of every input in every state can then be written as a table, with one row per state
and one column per input, and the machine can be tested exhaustively by visiting every cell of that
table. Placed before an actuator, such a machine is a gate: it admits a command only if the system's
current state can accept it, whatever produced the command. \TODO{a source for finite-state
machines in control software, e.g.\ Harel, ``Statecharts'', Science of Computer Programming 1987 ---
not in references.bib}

\paragraph{Publish and subscribe.} A command bus decouples the stages that produce commands from the
stage that consumes them. Producers \emph{publish} messages to the bus without addressing any
receiver, and consumers receive them without addressing any sender, so either side can be replaced
without changing the other. A message may still name its origin in its content. \TODO{a survey of
publish/subscribe, e.g.\ Eugster et al., ``The many faces of publish/subscribe'', ACM Computing
Surveys 2003 --- not in references.bib} A message is commonly serialised as text, for instance as a
\gls{json} object, and carried over a network protocol. The \gls{udp} carries each message as one
self-contained datagram, with no connection to set up and no retransmission. It adds little delay,
but it neither guarantees delivery nor preserves order.

\paragraph{Ordering.} When two producers publish to one consumer, the order in which messages
arrive need not be the order in which they were issued. A message issued first can arrive second
because its producer was slower, or because the transport reordered it. A shared counter that
issues each message a \emph{sequence number} restores a single order. Every message carries its
number, and the consumer can compare the numbers of any two messages to tell which was issued
first. A consumer can then discard a message made stale by one issued after it, whatever their
order of arrival.

\section{Foundations defined in the \emph{M\'emoire de Master}}
\label{sec:bg-master}
The language model on which the parse of a spoken command depends is the subject of the
\emph{M\'emoire de Master}, whose Background chapter defines its foundations. They are named here in
one sentence each and not repeated. A language model assigns probabilities to token sequences and
generates text one token at a time with a transformer decoder~\cite{vaswani2017}, and an \gls{slm} is
such a model at the small end of the scale. \Gls{lora} fits such a model to a task by training a
small number of added parameters while its original weights stay frozen~\cite{lora}. Post-training
quantisation stores the weights at reduced precision, and the \gls{gguf}\glsadd{ggml} formats Q8\_0 and Q4\_K\_M
of \texttt{llama.cpp}~\cite{llamacpp} are the two levels that document compares. Grammar-constrained
decoding under a \gls{gbnf} grammar masks, at each step, every token that would take the output
outside a formal language. Every output is then a string of that language or, if cut off at a token
limit, a prefix of one. A spoken command, once transcribed, is parsed into an intent and a set of
slot values, the frame of spoken-language understanding~\cite{qin2021}. A transcript is scored by
the \gls{wer}, the minimum number of word edits that turn it into its reference, divided by the
number of reference words. The parser is scored by \gls{em}, by intent and slot F1, by schema
validity, and by the false-command and safe-failure rates, and two parsers are compared on the same
items by McNemar's paired test~\cite{mcnemar1947}.

\section{Conclusion}
\label{sec:bg-conclusion}
This chapter has defined the stages of a speech pipeline, from the keyword spotter and its
operating point to endpointing, recognition and the \gls{snr} at which it is tested. It has
described the quadrotor, the nested loops that fly it and the \gls{pid} law they use, and the
three elements of swarm control: flocking, formations with their slot assignment, and separation
by potential fields, together with the simulators on which a controller is tested. It has set out
the timing vocabulary of a real-time system: latency percentiles, budgets that attribute a miss,
periodic work and preemption. It has described how four cores are shared under a general-purpose
operating system, and the state machine and command bus that carry and gate a command.
Chapter~\ref{chap:state-of-the-art} reviews the systems that published work has built from these
elements, compares them to one another, and identifies what none of them provides.
